%-----------------------------------------------
%                   Template
%-----------------------------------------------
%                   Settings
%-----------------------------------------------
%   Ändra språk, och snabb compile
\usepackage[english]{babel}       % Språk. [swedish] eller [english].
\usepackage[]{graphicx}           % Bilder. [demo] för snabb compile!

%-----------------------------------------------
%                   Packages
%-----------------------------------------------
\usepackage[utf8]{inputenc}          % Kodning av text.

\pagenumbering{gobble}       % Stoppar sidnumrering på titelsida.
\usepackage{mathtools}       % Ekvationer.
\usepackage{bm}              % Bold matte
\usepackage[a4paper,left=40mm,right=40mm,top=40mm,bottom=60mm]{geometry}      % Sidlayout.
\usepackage{float}           % För [H].
\usepackage[T1]{fontenc}     % Bra för åäö.
\usepackage{csquotes}        % Fancy Quotes
\usepackage{microtype}       % Fixar bättre textlayout.
\usepackage{amsmath}
\usepackage{tabto}
\usepackage{glossaries}
\usepackage{lipsum}

\usepackage[margin=15pt,font={small},labelfont={bf}%
,justification=centering]{caption}      % Snygg figurtext.
\usepackage{booktabs}                   % Snygga Tabeller.
\usepackage{longtable}                  % Långa tabeller, flera sidor.
\setlength{\abovecaptionskip}{3pt}      % Flyttar tabelltext närmare.
\renewcommand{\arraystretch}{1.2}       % Gör tabeller  större.
\setlength{\parindent}{0cm}             % Tar bort "indent" på ny rad.
%\usepackage[none]{hyphenat}             % Slutar dela upp ord.

\usepackage[backend=biber,style=ieee]{biblatex}
\addbibresource{Essential/ref.bib} % Mapp för Källor.
\graphicspath{{Images/}}                % Mapp för bilder.
\usepackage{tikz}
\usetikzlibrary{external}
\tikzexternalize[prefix=tikz/]
\usepackage{pdfpages}                   % PDFs.
\addto{\captionsswedish}{\renewcommand*%%
{\contentsname}{Innehållsförteckning}}  % Ändrar namn -> Innehållsförteckning

% URLs
\usepackage{url}
\usepackage{hyperref}                   % Hyperlänkar + färg.
\usepackage[nameinlink,noabbrev,swedish]{cleveref} % \cref{} & \Cref{}.

% Allow URL breaks at / and -
\def\UrlBreaks{\do\/\do-}

% Hyperref settings
\usepackage{color}                                 % Färg.
\definecolor{royalblue}{rgb}{0.0, 0.14, 0.4}       % Egen färg.
\hypersetup{                                       % Färg på hyperlänkar.
     colorlinks  = true,
     linkcolor   = black,              % internal links.
     citecolor   = red,          % bibliography.
     filecolor   = royalblue,          % file links.
     urlcolor    = royalblue           % external links.
}

\usepackage{matlab-prettifier}          % Matlab text!
\usepackage[per-mode=symbol]{siunitx}   % si. Skriva enheter/nummer.
\sisetup{output-decimal-marker={,},%    % si. ("{,}" till "{.}" för engelska).
range-phrase=--,range-units=single,exponent-product=\cdot}

%-----------------------------------------------
%                 New commands                
%-----------------------------------------------
% Använd för radbyte. \n
\newcommand{\n}{\vskip 1em}

% Bold matte och text (Automatisk) ex. \B{TEXT och } $\B{\alpha}$
\newcommand*{\B}[1]{\ifmmode\bm{#1}\else\textbf{#1}\fi}

% Partial Derivative. \pde{a}{b}
\newcommand{\pde}[2]{\frac{\partial #1}{\partial #2}}

% Använd för matlab fil.
\newcommand{\matlabfile}[1]
{\UseRawInputEncoding\hyphenpenalty=50\exhyphenpenalty=50\lstinputlisting[style = Matlab-editor,basicstyle = \mlttfamily]{#1}}

% Använd för matlab text.
\lstnewenvironment{matlab}
{\UseRawInputEncoding\hyphenpenalty=50\exhyphenpenalty=50\lstset{style = Matlab-editor,basicstyle = \mlttfamily}}{}

% Använd för centrerade tabeller
\newcolumntype{P}[1]{>{\centering\arraybackslash}p{#1}}

%%%%%%%%%%%%%%%%%%%%%%%%%%%%%%%%%%%%%%%%%%%%%%%%%%%%
%%  Creative Commons CC BY 4.0, Hugo Laestander %%%%
%%%%%%%%%%%%%%%%%%%%%%%%%%%%%%%%%%%%%%%%%%%%%%%%%%%%